\documentclass[11pt]{article}
\usepackage[utf8]{inputenc}
\usepackage[spanish]{babel}
\usepackage{amsmath,amssymb,amsthm}
\usepackage{graphicx}
\usepackage[margin=2.5cm]{geometry}
\usepackage{enumitem}

\theoremstyle{plain}
\newtheorem{teorema}{Teorema}
\theoremstyle{definition}
\newtheorem{definicion}{Definición}
\newtheorem{ejemplo}{Ejemplo}

\title{Semana 4 -- Materia}
\author{Tutorías -- Introducción al Cálculo}
\date{24/08--28/08/2026}

\begin{document}
\maketitle

\section{Intervalos (sección 2.4 del libro)}

La sección 2.4, ``Conjuntos de Números Reales'', define formalmente los subconjuntos de
$\mathbb{R}$ más usados a partir de los axiomas de existencia de conjuntos (Axiomas 2.17 y 2.18).
Entre ellos están los \textbf{intervalos}, que son los conjuntos que usaremos constantemente al
escribir conjuntos solución de inecuaciones.

\subsection{Definiciones}

Dados $a,b \in \mathbb{R}$ con $a<b$:
\begin{itemize}[nosep]
  \item Intervalo cerrado: $[a,b] = \{x\in\mathbb{R} : a\le x \le b\}$.
  \item Intervalo abierto: $\,]a,b[\, = \{x\in\mathbb{R} : a<x<b\}$.
  \item Cerrado en $a$, abierto en $b$: $[a,b[\, = \{x\in\mathbb{R} : a\le x<b\}$.
  \item Abierto en $a$, cerrado en $b$: $\,]a,b] = \{x\in\mathbb{R} : a<x\le b\}$.
  \item Infinito por la izquierda, abierto en $b$: $\,]-\infty,b[\, = \{x\in\mathbb{R} : x<b\}$.
  \item Infinito por la izquierda, cerrado en $b$: $\,]-\infty,b] = \{x\in\mathbb{R} : x\le b\}$.
  \item Infinito por la derecha, abierto en $a$: $\,]a,\infty[\, = \{x\in\mathbb{R} : a<x\}$.
  \item Infinito por la derecha, cerrado en $a$: $[a,\infty[\, = \{x\in\mathbb{R} : a\le x\}$.
\end{itemize}
Además $\,]-\infty,\infty[\,$ se usa para denotar a $\mathbb{R}$ completo. Note que el libro usa la
notación ``a la francesa'' con corchetes invertidos ($]a,b[$) para los intervalos abiertos, en vez de
paréntesis; mantenemos esa notación por consistencia con el texto.

Sobre conjuntos en general (no solo intervalos) se pueden aplicar las operaciones usuales: unión
$A\cup B$, intersección $A\cap B$, diferencia $A-B$ y la relación de subconjunto $A\subseteq B$
(Definiciones 2.22--2.25), todas definidas igual que en teoría de conjuntos elemental.

\begin{ejemplo}
$[-2,1] \cap \,]0,15] = \,]0,1]$.
\end{ejemplo}

\begin{ejemplo}
$\big([1,3[\, \cup \,]4,6]\big) \cap \,]2,5[\, = \,]2,3[\, \cup \,]4,5[$.
\end{ejemplo}

\begin{ejemplo}
Escriba $\mathbb{R} - \{2\}$ como unión de intervalos.

\emph{Solución.} $\mathbb{R}-\{2\} = \,]-\infty,2[\, \cup \,]2,\infty[$.
\end{ejemplo}

\section{Inecuaciones lineales y cuadráticas (sección 2.6)}

\subsection{Ecuaciones de referencia}

Antes de ver inecuaciones conviene recordar (Teorema 2.5) cómo se resuelven las ecuaciones
asociadas. La \textbf{ecuación de primer grado} $ax+b=0$ (con $a\neq 0$) tiene siempre una única
solución $x=-b/a$. La \textbf{ecuación de segundo grado} $ax^2+bx+c=0$ (con $a\neq 0$), llamando
$\Delta = b^2-4ac$:
\begin{itemize}[nosep]
  \item Si $\Delta \ge 0$, las soluciones son $x_1 = \dfrac{-b+\sqrt{\Delta}}{2a}$ y
  $x_2 = \dfrac{-b-\sqrt{\Delta}}{2a}$, y se tiene la factorización $ax^2+bx+c = a(x-x_1)(x-x_2)$.
  \item Si $\Delta < 0$, la ecuación no tiene soluciones reales.
\end{itemize}

\subsection{Inecuación de primer grado (2.6.5)}

Una inecuación $ax+b>0$ (con $a\neq 0$; análogamente con $\ge$, $<$, $\le$) se resuelve despejando
$x$ directamente: como $a\neq 0$,
\[
  ax+b>0 \;\leftrightarrow\; x>-\frac{b}{a} \quad (\text{si } a>0).
\]
\textbf{Cuidado:} si $a<0$, al dividir por $a$ la desigualdad se invierte (Teorema 2.3 (XI)). El
conjunto solución de $ax+b>0$ es $\,]-b/a,\infty[$ cuando $a>0$. Reemplazando $>$ por $\ge$, $<$ o
$\le$ se obtienen, análogamente, $[-b/a,\infty[$, $\,]-\infty,-b/a[$ y $\,]-\infty,-b/a]$
(para $a>0$; si $a<0$ los intervalos se invierten).

\begin{ejemplo}
Resolver $3x-6>0$.

\emph{Solución.} $3x>6 \leftrightarrow x>2$. Conjunto solución: $\,]2,\infty[$.
\end{ejemplo}

\begin{ejemplo}
Resolver $-2x+4 \ge 0$.

\emph{Solución.} $-2x \ge -4$. Como el coeficiente de $x$ es negativo, al despejar se invierte la
desigualdad: $x \le 2$. Conjunto solución: $\,]-\infty,2]$.
\end{ejemplo}

\subsection{Inecuación de segundo grado (2.6.6)}

Para $ax^2+bx+c>0$ (con $a\neq 0$), el signo de la expresión depende del discriminante
$\Delta=b^2-4ac$. El teorema correspondiente (sección 2.6.6) dice:

\begin{teorema}[Signo de la cuadrática]
Sean $a,b,c\in\mathbb{R}$, $\Delta=b^2-4ac$, $a\neq 0$.
\begin{enumerate}[label=(\Roman*)]
  \item Si $\Delta \ge 0$, sean $x_1\le x_2$ las soluciones de $ax^2+bx+c=0$:
  \begin{enumerate}[label=(\alph*)]
    \item Si $a>0$: $ax^2+bx+c>0 \leftrightarrow (x>x_2 \lor x<x_1)$.
    \item Si $a<0$: $ax^2+bx+c>0 \leftrightarrow x_1<x<x_2$.
  \end{enumerate}
  \item Si $\Delta<0$:
  \begin{enumerate}[label=(\alph*)]
    \item Si $a>0$: $\forall x\in\mathbb{R}\ (ax^2+bx+c>0)$.
    \item Si $a<0$: $\forall x\in\mathbb{R}\ (ax^2+bx+c<0)$.
  \end{enumerate}
\end{enumerate}
\end{teorema}

La demostración (caso $\Delta\ge0$, $a>0$) parte de la factorización $ax^2+bx+c=a(x-x_1)(x-x_2)$ y
usa que un producto de dos factores es positivo si y solo si ambos factores tienen el mismo signo:
\[
  a(x-x_1)(x-x_2)>0 \;\leftrightarrow\; \big((x-x_1)>0 \land (x-x_2)>0\big) \lor \big((x-x_1)<0 \land (x-x_2)<0\big)
  \;\leftrightarrow\; x>x_2 \lor x<x_1.
\]
Este mismo razonamiento ``signo del producto de factores'' es la técnica general que usamos para
armar \textbf{tablas de signos}: se factoriza la expresión, se ubican las raíces de cada factor en la
recta real, y se analiza el signo de cada factor (y por lo tanto el signo del producto) en cada
intervalo determinado por esas raíces.

\begin{ejemplo}
Resolver $x^2-x-6>0$.

\emph{Solución.} $\Delta = (-1)^2-4(1)(-6)=25>0$, $x_1=-2$, $x_2=3$ (raíces de $x^2-x-6=0$), y
$a=1>0$. Factorizando: $x^2-x-6=(x-3)(x+2)$. Tabla de signos:
\[
\begin{array}{c|ccccc}
  & x<-2 & x=-2 & -2<x<3 & x=3 & x>3 \\\hline
  (x+2) & - & 0 & + & + & + \\
  (x-3) & - & - & - & 0 & + \\\hline
  (x+2)(x-3) & + & 0 & - & 0 & +
\end{array}
\]
Como buscamos $(x+2)(x-3)>0$, el conjunto solución es $\,]-\infty,-2[\, \cup \,]3,\infty[$, lo que
coincide con el caso (I)(a) del teorema ($a>0$: solución fuera de las raíces).
\end{ejemplo}

\begin{ejemplo}
Resolver $-x^2+2x+3 \ge 0$.

\emph{Solución.} Aquí $a=-1<0$. $\Delta=4+12=16>0$, raíces $x_1=-1$, $x_2=3$. Por el caso (I)(b)
del teorema (con $\ge$ en vez de $>$, incluyendo los extremos): el conjunto solución es $[-1,3]$.
\end{ejemplo}

\begin{ejemplo}
Resolver $x^2+x+1>0$.

\emph{Solución.} $\Delta = 1-4=-3<0$ y $a=1>0$. Por el caso (II)(a) del teorema, la desigualdad se
cumple para todo $x\in\mathbb{R}$. Conjunto solución: $\mathbb{R}$.
\end{ejemplo}

\section{Inecuaciones racionales y por factores (Problema 2.3)}

El libro no dedica todavía una sección separada a ``inecuaciones racionales''; el primer ejemplo de
este tipo aparece como un problema resuelto (Problema 2.3) que combina un factor cuadrático con
raíces reales y otro factor cuadrático \emph{sin} raíces reales (es decir, de signo constante). La
idea clave es: si uno de los factores de un producto nunca cambia de signo, se puede ``eliminar''
ese factor sin cambiar el sentido de la desigualdad (dividiendo por él, ya que es siempre positivo o
siempre negativo).

\begin{ejemplo}[Problema 2.3 del libro]
Resolver $(x^2+3x-4)(2x^2+4)>0$.

\emph{Solución.} La ecuación $2x^2+4=0$ no tiene soluciones reales (su discriminante es
$0^2-4\cdot2\cdot4=-32<0$) y el coeficiente principal es $2>0$, luego por el Teorema de la sección
2.6.6, $2x^2+4>0$ para todo $x\in\mathbb{R}$. Como este factor siempre es positivo, multiplicar o
dividir por él no cambia el sentido de la desigualdad:
\[
  (x^2+3x-4)(2x^2+4)>0 \;\leftrightarrow\; x^2+3x-4>0 \;\leftrightarrow\; (x+4)(x-1)>0
  \;\leftrightarrow\; (x>1 \lor x<-4).
\]
Conjunto solución: $\,]-\infty,-4[\, \cup \,]1,\infty[$.
\end{ejemplo}

Esta misma idea (tabla de signos con todos los factores, incluyendo los que van en el denominador)
es la que se usará la próxima semana para inecuaciones racionales más generales, donde además hay
que excluir del conjunto solución los valores que anulan el denominador.

\end{document}
